\documentclass[12pt]{article}
\usepackage[margin=1in]{geometry}
\usepackage{enumitem}
\usepackage{titlesec}
\usepackage{parskip}
\usepackage{graphicx}
\usepackage{xcolor}
\usepackage{hyperref}
\usepackage{fontenc}

\title{\textbf{Strom 1990 RG Notes}\\
\large Strom, Kaare. ``A Behavioral Theory of Competitive Political Parties'' [1990].\\
\textit{American Journal of Political Science} 34(2): 565.}
\author{}
\date{}

\begin{document}

\maketitle

Builds on ``policy, office, or votes'' model to fill in the ``gaps'' -- what about when none of these three models can appropriately explain party behavior?

Vote-seeking: Downs

Office-seeking: coalition theory, Riker

Critique of vote-seeking: simply does not align with contemporary politics particularly well; very large number of exceptions etc.

Critique of office-seeking: not supported by behavior of minority coalition members, and the general apathy of some parties to be involved in a coalition government

Critique of policy: not supported by broad delegation to lower political bodies

This article moves towards two ``breaks'' of current assumptions'': begins to treat parties as constrained by political agreements, and begins to evaluate the supply of political goods rather than just the demand (that is, the provision of spoils)

Create a 3-space representation of behavior over three dimensions: policy, office, and votes, whereby parties have some assignment of relative weights between those outcomes.

Party leaders may be entrepreneurial and seek to gain office particularly to extract rents, but also must establish and maintain rent-costly institutions to actually manage and organize the party as a whole

\begin{center}
\includegraphics[width=0.85\textwidth]{figure-1.png}
\end{center}

\begin{center}
\includegraphics[width=0.85\textwidth]{figure-2.png}
\end{center}

See Harmel \& Janda 1994 discussing different goals of parties. Also note that parties can be ``democracy-seeking''

\end{document}
